\section{Устойчивость и метод Ляпунова}
\label{sec:stability}

\subsection{Что именно требуется доказать об устойчивости}
\label{subsec:stability-definitions}
\noindent\textit{Источник: L3, слайды 1--10.}

Лекция посвящена анализу устойчивости с помощью функции Ляпунова.
Для линейных систем в ней перечислены критерии Рауса--Гурвица,
Найквиста и Михайлова; для нелинейных --- круговой критерий,
критерий Попова и метод функций Ляпунова. Сами частотные критерии
в L3 не выводятся. Метод А. М. Ляпунова позволяет исследовать
устойчивость без явного решения дифференциальных уравнений.

Рассмотрим систему и её равновесие
\[
 \dot x=f(x),\qquad x(t_0)=x_0\in\mathbb R^n,\qquad f(0)=0.
\]
Здесь $x^*=0$ --- состояние, которое должно сохраняться при
отсутствии отклонений. Определения отвечают на разные вопросы.
\begin{enumerate}[leftmargin=*,itemsep=3pt]
 \item \textbf{Устойчивость по Ляпунову:} достаточно малое начальное
 отклонение остаётся малым при всех последующих временах:
 \[
  \forall\varepsilon>0\ \forall t_0\ge0\ \exists\delta(\varepsilon,t_0)>0:
  \quad \norm{x_0}<\delta\ \Longrightarrow\
  \norm{x(t)}<\varepsilon\quad\forall t\ge t_0.
 \]
 Сходимость к нулю здесь не требуется. На рисунке слайда 6
 траектория остаётся внутри трубки радиуса $\varepsilon$, хотя
 может колебаться и не приближаться к оси.
 \item \textbf{Асимптотическая устойчивость:} устойчивость по Ляпунову
 дополняется притяжением некоторой окрестности:
 \[
  \exists\delta_a>0:\quad \norm{x_0}<\delta_a
  \ \Longrightarrow\ \lim_{t\to\infty}x(t)=0.
 \]
 Поэтому одного наблюдения $x(t)\to0$ недостаточно: требуется
 также устойчивость.
 \item \textbf{Равномерность} означает независимость соответствующих
 оценок от начального момента $t_0$. \textbf{Глобальность}
 означает, что область притяжения охватывает всё пространство
 начальных состояний. Это разные свойства: первое относится
 ко времени начала движения, второе --- к величине $x_0$.
 \item \textbf{Экспоненциальная устойчивость} задаёт скорость
 затухания: существуют $c\ge1$, $\sigma>0$ и окрестность нуля, где
 \[
  \norm{x(t)}\le c\,e^{-\sigma(t-t_0)}\norm{x_0},\qquad t\ge t_0.
 \]
 При выполнении оценки для всех $x_0\in\mathbb R^n$ устойчивость
 глобально экспоненциальная.
 \item \textbf{Устойчивость за конечное время (finite-time):}
 равновесие устойчиво и для каждого допустимого $x_0$ существует
 конечное время установления $T(x_0)$, после которого
 \[
  x(t)=0\quad\forall t\ge t_0+T(x_0).
 \]
 \item \textbf{Устойчивость за фиксированное время (fixed-time):}
 дополнительно существует общая граница
 \[
  T(x_0)\le T_{\max}<\infty\quad\forall x_0\in\mathbb R^n.
 \]
 В первом случае время может неограниченно расти с $\norm{x_0}$,
 во втором одна оценка действует для всех начальных состояний.
\end{enumerate}

\textbf{Дополнение для защиты.}
Для автономной системы сдвиг $t_0$ не меняет уравнение, поэтому
обычные локальные определения устойчивости и асимптотической
устойчивости дают равномерность по $t_0$.
На слайде 8 экспоненциальная оценка записана сокращённо как
$\norm{x(t)}\le ce^{-\sigma t}$; в стандартной равномерной оценке
необходимы множитель $\norm{x_0}$ и время $t-t_0$, как выше.
Для равновесия $x^*\ne0$ все определения применяются к
отклонению $x-x^*$.

\subsection{Прямой метод Ляпунова: условия и геометрический смысл}
\label{subsec:stability-direct}
\noindent\textit{Источник: L3, слайды 2, 11--17.}

Пусть $D\subset\mathbb R^n$ --- открытая область, содержащая нуль,
а $f:D\to\mathbb R^n$ локально липшицева.
Это означает, что около каждой точки существует конечная
константа $L$, для которой
\[
 \norm{f(x)-f(y)}\le L\norm{x-y}.
\]
График на слайде 12 иллюстрирует ограниченность локального наклона:
не требуется ни глобальная константа $L$ на всём пространстве,
ни линейность $f$.

Ищем непрерывно дифференцируемую скалярную функцию
$V:D\to\mathbb R$, положительно определённую относительно нуля:
\[
 V(0)=0,\qquad V(x)>0\quad\text{при }x\in D\setminus\{0\}.
\]
Она играет роль обобщённой энергии. Её производная вычисляется
\emph{вдоль решений рассматриваемой системы}:
\[
 \dot V(x)=\sum_{i=1}^{n}\frac{\partial V}{\partial x_i}\dot x_i
 =\sum_{i=1}^{n}\frac{\partial V}{\partial x_i}f_i(x)
 =\frac{\partial V}{\partial x}f(x).
\]
В лекции $\partial V/\partial x$ --- строка, то есть градиент
в принятой там записи. После выбора управления в эту формулу
нужно подставлять правую часть \emph{замкнутой} системы.

\textbf{Теорема Ляпунова.}
Если $\dot V(x)\le0$ в $D$, то нуль устойчив по Ляпунову.
Если $\dot V(x)<0$ для всех $x\in D\setminus\{0\}$,
то нуль асимптотически устойчив.
На слайде 16 условие строгой отрицательности следует понимать
именно вне нуля: в равновесии $\dot V(0)=0$.

Чаша на слайдах 2 и 17 изображает график $V$.
Градиент направлен к возрастанию её высоты.
Условие $\dot V\le0$ означает, что вектор скорости $f(x)$
направлен внутрь подуровней $\{x:V(x)\le c\}$ или касается их
границы. Оно не требует, чтобы $f(x)$ совпадал с
$-\operatorname{grad}V$: траектория может вращаться, одновременно
теряя энергию.

\textbf{Дополнение для защиты: почему работает теорема.}
На достаточно малой сфере $\norm{x}=\varepsilon$ функция $V$
имеет положительный минимум $m_\varepsilon$.
Из непрерывности $V(0)=0$ можно выбрать $\delta$, чтобы
$\norm{x_0}<\delta$ влекло $V(x_0)<m_\varepsilon$.
Поскольку $V$ не возрастает, траектория не пересечёт эту сферу.
Для строгого знака на любом компактном кольце
$r\le\norm{x}\le\varepsilon$ имеем $\dot V\le-d_r<0$.
Траектория не может оставаться в таком кольце неограниченно долго:
иначе положительная $V$ стала бы отрицательной.
Это даёт притяжение к нулю.

Для \emph{глобального} вывода недостаточно показать знак производной
в небольшой окрестности. Достаточны $D=\mathbb R^n$,
глобальная положительная определённость и радиальная
неограниченность
\[
 \norm{x}\to\infty\quad\Longrightarrow\quad V(x)\to\infty.
\]
Тогда подуровни компактны: убывание $V$ ограничивает всё состояние
и исключает уход решения на бесконечность за конечное время.
Если, кроме того,
\[
 a_1\norm{x}^2\le V(x)\le a_2\norm{x}^2,
 \qquad \dot V(x)\le-a_3\norm{x}^2,\qquad a_i>0,
\]
то интегрирование $\dot V\le-(a_3/a_2)V$ даёт
\[
 \norm{x(t)}\le\sqrt{\frac{a_2}{a_1}}\,
 e^{-\frac{a_3}{2a_2}(t-t_0)}\norm{x_0}.
\]
Следовательно, отрицательный знак доказывает сходимость,
а количественная оценка производной позволяет оценить её скорость.

\subsection{Квадратичная функция и связь с линеаризацией}
\label{subsec:stability-linearization}
\noindent\textit{Источник: L3, слайды 18--21.}

Около равновесия гладкую систему представляют как
\[
 \dot x=Ax+g(x),\qquad
 A=\left.\frac{\partial f}{\partial x}\right|_{x=0},
 \qquad g(x)=o(\norm{x}).
\]
Здесь $g(x)$ --- нелинейный остаток; его не следует путать с
коэффициентом при управлении из других лекций.
Выберем $V(x)=x^TPx$, где $P=P^T\succ0$. Тогда
\[
 \begin{aligned}
 \dot V
 &=x^T(A^TP+PA)x+2x^TPg(x),\\
 A^TP+PA&=-Q,\qquad Q=Q^T\succ0,\\
 \dot V&=-x^TQx+2x^TPg(x).
 \end{aligned}
\]
Если в рассматриваемой области $\norm{g(x)}\le\gamma\norm{x}$,
то для евклидовой нормы и согласованной матричной нормы
\[
 \dot V\le-\bigl[\lambda_{\min}(Q)-2\gamma\norm{P}\bigr]\norm{x}^2.
\]
Поэтому достаточно
\[
 \gamma<\frac{\lambda_{\min}(Q)}{2\norm{P}}.
\]
Смысл условия: затухание линейной части должно преобладать
над нелинейным остатком.

\textbf{Дополнение для защиты.}
Если $A$ гурвицева, то для каждого $Q\succ0$ существует
единственное $P\succ0$, решающее уравнение Ляпунова:
\[
 P=\int_0^\infty e^{A^Tt}Qe^{At}\,dt.
\]
Интеграл сходится из-за экспоненциального убывания $e^{At}$;
интегрирование производной $e^{A^Tt}Qe^{At}$ даёт
$A^TP+PA=-Q$.
Условие $g=o(\norm{x})$ позволяет сделать $\gamma$ сколь угодно
малой выбором окрестности. Так объясняется локальный вывод
об устойчивости нелинейной системы по гурвицевой линеаризации.
Глобальный вывод требует глобальных оценок.

\textbf{Пример 1 из лекции.}
Для $\dot x=ax^3$ матрица линеаризации равна
$A=f'(0)=0$, независимо от $a$. Однако при $V=x^2$
\[
 \dot V=2ax^4.
\]
При $a=0$ все состояния постоянны: нуль устойчив, но не притягивает.
При $a<0$ нуль глобально асимптотически устойчив.
При $a>0$ ненулевые состояния движутся от нуля, и он неустойчив.
Таким образом, нулевое собственное значение линеаризации
само по себе не определяет устойчивость.

\textbf{Дополнение для защиты.}
При $a<0$ решение имеет вид
\[
 x(t)=\frac{x_0}{\sqrt{1+2|a|x_0^2(t-t_0)}}.
\]
Оно стремится к нулю алгебраически; это пример асимптотической
устойчивости без экспоненциальной.
В работе 3 устойчивость нуль-динамики нужно исследовать так же
внимательно: нулевой корень не даёт асимптотического затухания
автоматически.

\subsection{Критерии сходимости за конечное и фиксированное время}
\label{subsec:stability-finite-fixed}
\noindent\textit{Источник: L3, слайды 9--10, 22--23.}

Пусть $V$ непрерывно дифференцируема и положительно определена.
Если вдоль траекторий
\[
 \dot V(x)\le-cV^\alpha(x),\qquad c>0,\quad0<\alpha<1,
\]
то нуль устойчив за конечное время, причём
\[
 T(x_0)\le\frac{V^{1-\alpha}(x_0)}{c(1-\alpha)}.
\]
Если вместо этого
\[
 \dot V(x)\le-c_1V^\alpha(x)-c_2V^\beta(x),\qquad
 c_1,c_2>0,\quad0<\alpha<1<\beta,
\]
то получаем фиксированную границу
\[
 T_{\max}\le\frac{1}{c_1(1-\alpha)}
                +\frac{1}{c_2(\beta-1)}.
\]
Для глобальных утверждений условия должны действовать глобально,
а $V$ должна ограничивать всё состояние, например быть радиально
неограниченной.

\textbf{Дополнение для защиты: вывод оценок.}
Пока $V>0$,
\[
 \frac{d}{dt}V^{1-\alpha}
 =(1-\alpha)V^{-\alpha}\dot V\le-c(1-\alpha).
\]
Правая часть постоянна, поэтому положительная $V$ обязана
достичь нуля за указанное время. Для fixed-time сначала
рассматривают $V\ge1$: время спуска с любого $V(x_0)$ до единицы
не больше
\[
 \int_1^\infty\frac{dv}{c_2v^\beta}
 =\frac{1}{c_2(\beta-1)}.
\]
При $0<V\le1$ время до нуля не больше
$\int_0^1dv/(c_1v^\alpha)=1/[c_1(1-\alpha)]$.
Именно степень $\beta>1$ ограничивает время движения издалека,
а $\alpha<1$ обеспечивает конечное время около нуля.

Здесь нельзя автоматически переносить предположение о локальной
липшицевости $f$ в нуле из обычной теоремы Ляпунова.
Для локально липшицева автономного поля единственность решений
не позволяет ненулевой траектории слиться с равновесием за
конечное время. Finite-time-системы используют нелипшицевы
законы, например $\dot x=-|x|^p\sgn x$ при $0<p<1$,
либо разрывные законы с соответствующим понятием решения.
В работе 5 из условия $s\dot s\le-g_0\beta_0|s|$,
$g_0,\beta_0>0$, при $V=s^2/2$ следует
\[
 \dot V\le-g_0\beta_0\sqrt{2}\,V^{1/2},\qquad
 T\le\frac{|s(0)|}{g_0\beta_0}.
\]
Это конечное время достижения поверхности $s=0$.
Оно само по себе не означает достижения всего состояния нуля
за конечное время: движение вдоль поверхности исследуется отдельно.

\subsection{Примеры с маятником: энергия и недостаточность нестрогого знака}
\label{subsec:stability-pendulum}
\noindent\textit{Источник: L3, слайды 24--32.}

\textbf{Пример 2: маятник без трения.}
Пусть $x_1$ --- угол от нижней вертикали, $x_2=\dot x_1$,
$l>0$ --- длина, $m>0$ --- масса, $g>0$ --- ускорение тяжести.
На рисунке высота подъёма груза $h=l(1-\cos x_1)$,
его линейная скорость $v=lx_2$. Система и полная энергия:
\[
 \dot x_1=x_2,\qquad \dot x_2=-\frac{g}{l}\sin x_1,
 \qquad
 V=W_{\mathrm K}+W_{\mathrm П}
   =mgl(1-\cos x_1)+\frac12ml^2x_2^2.
\]
Около $(0,0)$ энергия положительно определена, а
\[
 \dot V=mgl\,x_2\sin x_1
       +ml^2x_2\left(-\frac{g}{l}\sin x_1\right)=0.
\]
Нуль устойчив. Ненулевое начальное значение энергии сохраняется,
поэтому траектория с такой энергией не может стремиться к нулю.
Колебания на слайде 27 не затухают: они идут по линиям постоянной
энергии, а не спускаются к её минимуму.

Линеаризация и её собственные значения:
\[
 A=\begin{pmatrix}0&1\\-g/l&0\end{pmatrix},\qquad
 \lambda_{1,2}=\pm j\sqrt{g/l}.
\]
Чисто мнимые корни не позволяют сделать вывод об асимптотической
устойчивости нелинейной системы по первому приближению.

\textbf{Дополнение для защиты: исправления и область вывода.}
На слайде 28 в собственных значениях пропущен квадратный корень:
должно быть $\pm j\sqrt{g/l}$, а не $\pm jg/l$.
Надпись $V(x_0)=x_0$ на слайде 27 следует читать как
$V(x(t))=V(x_0)$.
Функция $1-\cos x_1$ положительно определена относительно
$x_1=0$ только в выбранной окрестности: при $x_1=2\pi k$
она также обращается в нуль. Поэтому энергия не доказывает
глобальной устойчивости одного равновесия на всей плоскости
$(x_1,x_2)\in\mathbb R^2$.

\textbf{Пример 3: маятник с вязким трением.}
В принятой в лекции модели при $k>0$
\[
 \dot x_1=x_2,\qquad
 \dot x_2=-\frac{g}{l}\sin x_1-\frac{k}{m}x_2,
 \qquad \dot V=-l^2kx_2^2\le0.
\]
Энергия теряется за счёт скорости. Но её производная равна нулю
на всей прямой $x_2=0$, включая положения с $x_1\ne0$.
По прямой теореме уже доказана устойчивость, но строгий критерий
асимптотической устойчивости неприменим.
Рисунки слайда 31 подчёркивают: убывающая неотрицательная $V(t)$
может иметь ненулевой предел. Отсутствие $x_1$ в $\dot V$ требует
анализа движения на множестве нулевой диссипации.

На фазовом портрете слайда 32 траектории закручиваются к разным
нижним положениям $(2\pi k,0)$; верхние положения
$((2k+1)\pi,0)$ разделяют области притяжения.
Большая начальная скорость может перенести маятник через верхнее
положение до остановки в другой потенциальной яме.
Следовательно, на всей плоскости нельзя утверждать сходимость
каждого решения именно к $(0,0)$.

\textbf{Дополнение для защиты.}
Важно не буквальное присутствие всех переменных в формуле
$\dot V$, а её знак: отрицательная определённость означает
$\dot V(x)<0$ в \emph{каждой} ненулевой точке.
Например, $-x_1^2x_2^2$ содержит обе переменные, но также
не является отрицательно определённой функцией.

\subsection{Принцип инвариантности ЛаСалля}
\label{subsec:stability-lasalle}
\noindent\textit{Источник: L3, слайды 33--34.}

Множество $\Omega\subset D$ называется \emph{положительно
инвариантным}, если каждое решение, начавшееся в $\Omega$,
остаётся в нём при всех $t\ge t_0$.
Пусть $\Omega$ компактно и положительно инвариантно,
$V\in C^1(D)$, а $\dot V\le0$ в $\Omega$.
Обозначим
\[
 E=\{x\in\Omega:\dot V(x)=0\},\qquad
 M=\text{наибольшее инвариантное множество в }E.
\]
\textbf{Теорема ЛаСалля:} любое решение с $x_0\in\Omega$
приближается к $M$, то есть
$\operatorname{dist}(x(t),M)\to0$.
На схеме лекции $M\subset E\subset\Omega$:
попасть в $E$ и навсегда остаться в $E$ --- разные условия.

\textbf{Дополнение для защиты: идея доказательства.}
На компактном $\Omega$ функция $V$ ограничена снизу,
поэтому невозрастающее $V(x(t))$ имеет предел.
Предельные точки траектории образуют непустое инвариантное
множество, на котором $V$ постоянно и $\dot V=0$.
Оно содержится в $M$, откуда и следует приближение к $M$.
В отличие от прямой теоремы, для этой формулировки не требуется
положительная определённость $V$; для вывода об устойчивости
конкретного равновесия она по-прежнему полезна.

Для маятника с трением $E$ задаётся условием $x_2=0$.
Чтобы траектория оставалась на этой прямой, необходимо также
\[
 \dot x_2=-\frac{g}{l}\sin x_1=0.
\]
Следовательно, в $M$ остаются только равновесия $(k\pi,0)$.
В окрестности нижнего положения единственное из них --- $(0,0)$,
поэтому это равновесие асимптотически устойчиво.

\textbf{Дополнение для защиты: как выбрать правильное $\Omega$.}
Выбираем $0<c<2mgl$ и компоненту подуровня энергии около нуля:
\[
 \Omega_c=\{(x_1,x_2): |x_1|\le\pi,\ V(x_1,x_2)\le c\}.
\]
Это компактное положительно инвариантное множество.
На границах $x_1=\pm\pi$ энергия не меньше $2mgl>c$,
а пересечь границу $V=c$ наружу мешает $\dot V\le0$.
Внутри $\Omega_c$ из $x_2=0$, $\sin x_1=0$ следует $x_1=0$,
то есть $M=\{(0,0)\}$.
Одного условия $-\pi<x_1<\pi$ со слайда 34 недостаточно:
при большой скорости эта полоса не инвариантна.

\subsection{Устойчивость по входу--состоянию: определение и линейный случай}
\label{subsec:stability-iss-definition}
\noindent\textit{Источник: L3, слайды 35--39.}

Теперь рассматриваем систему с входом
\[
 \dot x=f(t,x,u),\qquad
 f:[0,\infty)\times\mathbb R^n\times\mathbb R^m\to\mathbb R^n.
\]
В лекции $f$ кусочно непрерывна по $t$, локально липшицева
по $x,u$, а $u(t)$ --- кусочно непрерывный ограниченный вход.
У невозмущённой системы $\dot x=f(t,x,0)$ предполагается
глобально равномерно асимптотически устойчивый нуль.
Следующий вопрос: насколько велико состояние при ненулевом входе?

Система называется \emph{устойчивой по входу--состоянию}
(input-to-state stable, ISS), если решения существуют для всех
$t\ge t_0$ и имеются функции $\beta\in\mathcal{KL}$,
$\gamma\in\mathcal K$, такие что
\[
 \norm{x(t)}\le
 \beta\bigl(\norm{x(t_0)},t-t_0\bigr)
 +\gamma\left(\sup_{t_0\le\tau\le t}\norm{u(\tau)}\right).
\]
Первое слагаемое затухает и описывает память о начальном состоянии.
Второе оценивает вклад входа. При $u=0$ получаем затухание к нулю;
ограниченный постоянный вход в общем случае оставляет ненулевое
отклонение.

\textbf{Дополнение для защиты: классы функций.}
Непрерывная $\gamma:[0,a)\to[0,\infty)$ принадлежит
$\mathcal K$, если она строго возрастает и $\gamma(0)=0$.
Если $a=\infty$ и $\gamma(r)\to\infty$ при $r\to\infty$,
это класс $\mathcal K_\infty$.
Функция $\beta(r,s)$ принадлежит $\mathcal{KL}$, если
$\beta(\cdot,s)\in\mathcal K$ при каждом фиксированном $s$,
а при фиксированном $r$ она убывает к нулю с ростом $s$.
Пример: $\beta(r,s)=ce^{-\sigma s}r$.

Для линейной системы $\dot x=Ax+Bu$ с гурвицевой $A$
формула вариации постоянных даёт
\[
 x(t)=e^{A(t-t_0)}x(t_0)
       +\int_{t_0}^{t}e^{A(t-\tau)}Bu(\tau)\,d\tau.
\]
При $\norm{e^{At}}\le ke^{-\lambda t}$, $k\ge1$, $\lambda>0$,
\[
 \begin{aligned}
 \norm{x(t)}
 &\le ke^{-\lambda(t-t_0)}\norm{x(t_0)}
  +\int_{t_0}^{t}ke^{-\lambda(t-\tau)}\norm{B}\norm{u(\tau)}\,d\tau\\
 &\le ke^{-\lambda(t-t_0)}\norm{x(t_0)}
  +\frac{k\norm{B}}{\lambda}
    \sup_{t_0\le\tau\le t}\norm{u(\tau)}.
 \end{aligned}
\]
Таким образом, $\beta(r,s)=ke^{-\lambda s}r$,
$\gamma(r)=k\norm{B}r/\lambda$; при $B=0$ можно взять
любую положительную линейную $\gamma\in\mathcal K$.

Идея нелинейного случая со слайда 38: взять функцию Ляпунова
невозмущённой системы и оценить добавку от входа.
Например, если
\[
 \norm{f(t,x,u)-f(t,x,0)}\le L\norm{u},
\]
то эта добавка мала при малом входе. Достаточно добиться
строгого убывания $V$ вне шара, радиус которого зависит
от величины входа. Состояние тогда притягивается к некоторой окрестности,
размер которой оценивается через уровни $V$.

\subsection{ISS-теорема Ляпунова, достаточная лемма и оба примера}
\label{subsec:stability-iss-lyapunov}
\noindent\textit{Источник: L3, слайды 40--43.}

\textbf{ISS-теорема.}
Пусть существует $V\in C^1([0,\infty)\times\mathbb R^n)$,
для которой
\[
 \alpha_1(\norm{x})\le V(t,x)\le\alpha_2(\norm{x}),
 \qquad\alpha_1,\alpha_2\in\mathcal K_\infty,
\]
и при $\norm{x}\ge\rho(\norm{u})$, $x\ne0$,
\[
 \frac{\partial V}{\partial t}
 +\frac{\partial V}{\partial x}f(t,x,u)
 \le-W_3(x),
\]
где $\rho\in\mathcal K$, а $W_3$ непрерывна и положительно
определена. Тогда система ISS, причём можно выбрать
\[
 \gamma=\alpha_1^{-1}\circ\alpha_2\circ\rho.
\]
Здесь $\circ$ означает композицию функций, а не их произведение.

\textbf{Дополнение для защиты: смысл оценки усиления.}
При $\norm{u}\le r$ производная отрицательна за пределами
шара $\norm{x}\le\rho(r)$.
В этом шаре $V\le\alpha_2(\rho(r))$; подуровень с такой высотой
лежит внутри шара радиуса
$\alpha_1^{-1}(\alpha_2(\rho(r)))$.
Так убывание обобщённой энергии превращается в оценку нормы
состояния. Условие отрицательности вне шара не утверждает,
что $\dot V$ отрицательна в каждой точке при ненулевом входе.

\textbf{Достаточная лемма.}
Если $f(t,x,u)$ непрерывно дифференцируема и глобально липшицева
по $(x,u)$ равномерно по $t$, а невозмущённый нуль глобально
экспоненциально устойчив, то система с входом ISS.
Это достаточное условие; неприменимость леммы не означает,
что ISS отсутствует.

\textbf{Дополнение для защиты: идея обоснования леммы.}
При этих предположениях обратная теорема Ляпунова даёт для
невозмущённой системы квадратичные оценки $V$ и оценки
$\dot V\le-a_3\norm{x}^2$,
$\norm{\partial V/\partial x}\le a_4\norm{x}$.
Глобальная липшицевость по входу позволяет записать
\[
 \dot V\le-a_3\norm{x}^2+a_4L\norm{x}\norm{u}
 \le-\frac{a_3}{2}\norm{x}^2
 \quad\text{при}\quad
 \norm{x}\ge\frac{2a_4L}{a_3}\norm{u}.
\]
Остаётся применить ISS-теорему.

\textbf{Пример 1: $\dot x=-x^3+u$.}
При $u=0$ нуль глобально асимптотически устойчив.
Возьмём $V=x^2/2$ и $0<\theta<1$. Тогда
\[
 \begin{aligned}
 \dot V&=-x^4+xu
       =-(1-\theta)x^4-\theta x^4+xu\\
       &\le-(1-\theta)x^4
       \quad\text{при}\quad
       |x|\ge\left(\frac{|u|}{\theta}\right)^{1/3}.
 \end{aligned}
\]
Действительно, при указанном условии
$|x||u|\le\theta|x|^4$.
Здесь $\alpha_1(r)=\alpha_2(r)=r^2/2$,
$\rho(r)=(r/\theta)^{1/3}$, поэтому
\[
 \gamma(r)=\left(\frac{r}{\theta}\right)^{1/3}.
\]
Этот пример показывает ISS без экспоненциального затухания
невозмущённой системы.

\textbf{Пример 2: $\dot x=-x-2x^3+(1+x^2)u^2$.}
Невозмущённая система глобально асимптотически устойчива,
но $f$ не глобально липшицева, поэтому предыдущая лемма
неприменима. Для $V=x^2/2$:
\[
 \dot V=-x^2-2x^4+x(1+x^2)u^2
       \le-x^4\quad\text{при }|x|\ge u^2.
\]
Проверка: $|x|u^2\le x^2$ и $|x|^3u^2\le x^4$
при $|x|\ge u^2$, поэтому положительная добавка
не больше $x^2+x^4$.
Имеем $\rho(r)=r^2$ и
\[
 \gamma(r)=r^2.
\]
Таким образом, непосредственная функция Ляпунова позволяет
доказать ISS за пределами условий достаточной леммы.

\subsection{Каскадные соединения и применение при защите работ}
\label{subsec:stability-cascade}
\noindent\textit{Источник: L3, слайд 44.}

Рассмотрим каскад
\[
 \dot x_1=f_1(t,x_1,x_2),\qquad
 \dot x_2=f_2(t,x_2).
\]
Здесь $x_1,x_2$ обозначают векторы подсистем.
Вторая подсистема не зависит от первой; её состояние выступает
входом первой.
Если первая подсистема ISS относительно входа $x_2$, а нуль
второй глобально равномерно асимптотически устойчив, то нуль
всего каскада глобально равномерно асимптотически устойчив
при обычных условиях существования решений.

\textbf{Дополнение для защиты: почему нужна именно ISS.}
Устойчивость второй подсистемы ограничивает $x_2(t)$ и даёт
$x_2(t)\to0$. ISS первой ограничивает $x_1(t)$ и исключает
конечный уход на бесконечность. Для сходимости применяем
ISS-оценку заново с момента $T$:
\[
 \norm{x_1(t)}\le
 \beta\bigl(\norm{x_1(T)},t-T\bigr)
 +\gamma\left(\sup_{T\le\tau\le t}\norm{x_2(\tau)}\right).
\]
Сначала выбираем $T$, чтобы весь оставшийся вход $x_2$ стал мал,
затем увеличиваем $t-T$, чтобы затухло первое слагаемое.
Получаем $x_1(t)\to0$. Совместная устойчивость следует из
устойчивости $x_2$ и ISS-оценки $x_1$ для малых начальных данных.
Так же объясняется свойство ISS: ограниченный вход, стремящийся
к нулю, вызывает состояние, стремящееся к нулю.

Просто асимптотической устойчивости двух подсистем при нулевых
входах недостаточно. Например,
\[
 \dot x_1=-x_1+x_1^2x_2,\qquad \dot x_2=-x_2
\]
имеет две устойчивые изолированные подсистемы при отключённой
связи, но при $x_1(0)>0$, $x_2(0)>0$,
$x_1(0)x_2(0)>2$ переменная $x_1$ уходит на бесконечность
за конечное время. Действительно,
\[
 \frac1{x_1(t)}=e^t\left[
 \frac1{x_1(0)}-\frac{x_2(0)}2(1-e^{-2t})\right]
\]
обращается в нуль. Здесь первая подсистема не ISS.

\textbf{Дополнение для защиты: связи с работами 3--5.}
В задаче глобальной стабилизации работы 3 после обратной связи
получен каскад
\[
 \dot\eta=-2\eta+\xi_1,\qquad
 \dot\xi_1=\xi_2,\qquad \dot\xi_2=-\xi_1-2\xi_2.
\]
Внешняя подсистема $(\xi_1,\xi_2)$ экспоненциально устойчива,
а внутреннее уравнение ISS относительно $\xi_1$:
\[
 |\eta(t)|\le e^{-2(t-t_0)}|\eta(t_0)|
 +\frac12\sup_{t_0\le\tau\le t}|\xi_1(\tau)|.
\]
По каскадной лемме все новые координаты сходятся к нулю.
Обратимый глобальный переход координат переносит этот вывод
на исходную систему.
В другой системе работы 3 нуль-динамика имеет собственные
значения $0,-2$: при
$V=(\eta_1^2+\eta_2^2)/2$ получаем
$\dot V=-(\eta_1-\eta_2)^2$ и целую прямую равновесий
$\eta_1=\eta_2$. Поэтому есть устойчивость, но нет
асимптотической устойчивости выбранной точки.
При неустойчивой нуль-динамике $\dot\eta=\eta$ затухание выхода
вообще не гарантирует ограниченности внутреннего состояния.

В работе 4 функция Ляпунова строится для всех ошибок
бэкстеппинга, например
$V_c=(x_1^2+z_2^2+z_3^2+z_4^2)/2$.
Глобальный вывод требует глобальной обратимости замены
и положительной определённости относительно нужного равновесия;
для третьей системы это $x^*=(0,1,0,0)$.
В работе 5 нужно различать конечное время попадания на поверхность,
устойчивость движения вдоль неё и эффект постоянного возмущения.
Нестрогое убывание энергии требует ЛаСалля либо другого
дополнительного аргумента, а ISS-оценка при ненулевом входе
обычно гарантирует окрестность равновесия, а не точное зануление.
